\chapter{Evolutor runtime above both engines}
\label{evod39:chapter}

Consider a workflow with five operations: retrieve a document containing
\emph{memory}, consume tokens $(1,2,1)$ in a native recurrence, consume the
same token IDs in a Transformer, add integers 7 and 8, and echo the sum through
a permitted tool. Predict which outputs can be compared directly. The sum is
15; the retrieved document is evidence; the two engine calls produce different
logit spaces. Neither engine's hidden state can be handed to the other.

This chapter builds Evolutor's coordination boundary above those distinct
computations. DOGMA is the non-Transformer direction; Hermon DNA is the
Transformer direction. The executable DOGMA-side fixture remains the
\emph{restricted Chapter 31 diagonal teaching recurrence}, not a complete
trained DOGMA model. The Hermon side imports Chapter 35, which serves the
actual Chapter 23 multi-layer Transformer with learned absolute positions.
These are executable continuity contracts, not evidence that the workflow
solves a learned genomic task.

\section{Make the workflow a typed object}
\label{evod39:typed}
\index{type checking}\index{single static assignment}
Evolutor owns the path from an objective to a plan, route, execution and
observable trace. It does not own the internal algebra of every backend.
The \Term{planner} decides which operations are needed and their dependencies.
The \Term{router} chooses an eligible implementation of one already specified
operation. A bad dependency graph and a missing backend are different errors.

\Plate{01-dataflow}{The executable fixture has typed values and separate control authority. Equal token IDs do not imply equal vocabulary semantics or engine outputs.}{evod39:dataflow}

Represent a value by $(\tau,x,\pi)$: type $\tau$, payload $x$, and provenance
$\pi$. Our finite alphabet contains integers, token sequences, logits,
documents and tool receipts. A token tuple alone is not a universal text
representation. The Chapter 31 fixture has vocabulary size three; Hermon's
causal ACGT head has size four. Deliberately using IDs in their intersection
permits a systems test but does not establish equivalent inputs as language
or biology. A real cross-model workflow needs an explicit tokenizer or
record conversion operation, with its own loss and coordinate contract.

A step is
\[
 s=(id,op,I,o,\tau_o,A,C),
\]
where $I$ names previously defined inputs, $o$ is its new output name, $A$
contains operation-specific attributes, and $C$ is a required engine capability
when appropriate. The plan pins a tenant/request pair, immutable input values,
an ordered step tuple and requested outputs. The tuple is a topological order
of a straight-line graph; branches, loops and concurrent effects are outside
this reference. Chapter 41 will make effects and lowering explicit in an IR.

\begin{center}
\begin{tabular}{lll}\toprule
Operation & Input type & Output type\\\midrule
engine & tokens & backend-specific logits\\
retrieve & version, query, immutable corpus attributes & documents\\
symbolic & integer, integer & integer\\
tool & integer, configured tool name & receipt\\\bottomrule
\end{tabular}
\end{center}
The source supplies \texttt{Step}, \texttt{Plan}, \texttt{Value} and
\texttt{validate}. Validation checks the entire graph before opening an
engine or calling a tool: unique IDs and SSA outputs, defined inputs, exact
operation signatures, required attributes and defined final outputs.
Integer validation excludes Python's Boolean subtype. A malformed later
operation cannot be discovered only after an earlier effect has run.

This is type preflight, not omniscience. A well-typed token can still exceed
one model's vocabulary or remaining context limit. The backend validates those
value-dependent constraints immediately before arithmetic. Likewise, an
earlier valid tool can have run before a later backend fails. We do not claim
that a plan is one atomic transaction.

\section{Route semantic requirements before optimizing cost}
\label{evod39:routing}
\index{routing!capability}\index{fallback}\index{cost model!routing}
For step $s$, let $C_s$ be required capabilities and $Q_s$ hard constraints.
Backend $e$ is eligible only when
\[
 C_s\subseteq C_e,\qquad Q_s(e)=\mathrm{true}.
\]
Artifact, device, precision, state continuity, tokenizer and policy constraints
belong in $Q_s$. Our minimal router requires one of
\texttt{compressed\_state} or \texttt{addressable\_kv} and an exact artifact
identity. It fails on zero matches or ambiguity. Multiple real eligible
replicas could use an explicit tie-break and cost policy; list order should
not silently choose semantics.
\Listing{Runtime.route}

The reference maps the former capability to the Chapter 31 adapter and the
latter to Chapter 35. This does not mean that every task has a unique natural
architecture. It means that an operation requiring continuation of a particular
native state cannot be satisfied by an unrelated backend merely because both
return logits. Fallback is legal only when the alternate implementation
preserves the same observable contract. Otherwise the planner must create
a new legal graph and record the semantic change.

After hard eligibility, a future policy may minimize
\[
 J(e)=w_L\widehat L_e+w_M\widehat M_e+w_C\widehat C_e+w_R\widehat R_e.
\]
Hats denote predictions: latency, memory, monetary or energy cost, and a
defined quality-risk proxy. Weights convert their units to a common score.
A large penalty is not a substitute for an authorization or compatibility
constraint. No timings or learned quality estimates are measured here;
Chapter 47 owns calibrated backend evidence.

\section{Keep complete state private and identity explicit}
\label{evod39:handles}
\index{state ownership}\index{Artifact binding}\index{continuation!chunk parity}
An \Term{opaque handle} contains adapter-instance identity, backend, artifact,
tenant/request and serial number. The adapter maps it to its private native
state. Equal shapes or equal numeric serials do not imply ownership.
\Plate{02-handles}{Handles address one live adapter mapping. Closing a handle invalidates that mapping; a different instance's serial zero remains distinct.}{evod39:handlesplate}

For Chapter 31, the sufficient state contains position $t$ and hidden vector
$h_t\in\mathbb R^D$, with fixed artifact parameters:
\[
 h_{t+1}=a\odot h_t+\sigma(w\odot E_{x_{t+1}})\odot E_{x_{t+1}},
 \qquad z_{t+1}=h_{t+1}^{\mathsf T}U.
\]
Chapter 35 instead holds each layer's ordered
$K^\ell,V^\ell\in\mathbb R^{H\times t\times d_h}$, consumed-token metadata,
causal mode and artifact binding. It computes learned position $P_t$ for the
next consumed token. There is no lawful conversion between these objects
implied by this shared handle envelope.

Shared method names such as \texttt{open}, \texttt{execute} and
\texttt{close} promise lifecycle behavior, not state representation.
A checkpoint operation could be added only with its native exact-continuation
contract: Chapter 33 has a concrete DOGMA checkpoint; a Hermon checkpoint
requires its complete KV and placement rules. This chapter does not invent a
universal checkpoint encoder or claim Chapter 36 page integration.

\Listing{EngineAdapter.execute}
Both imported routines validate before returning fresh candidate state.
The adapter publishes that candidate only after successful execution.
Invalid chunks leave the prior mapping unchanged. The returned logits are
immutable tuples, so callers cannot mutate the retained state through them.
The artifact is checked by the native routines, including in-place parameter
changes. Cleanup remains possible after an artifact mutation; otherwise an
error could permanently strand a live handle.

\begin{proposition}[Ownership and chunk refinement]
Assume serialized calls, trusted adapter code, fixed native artifacts and
non-mutating callers. A handle accepted by this adapter refers to exactly
one live state under its recorded owner. Splitting a valid token chunk into
nonempty pieces preserves the native concatenated output sequence in exact
arithmetic under the native transition contracts.
\end{proposition}
\begin{proof}
Every open uses a fresh serial within one adapter identity and records the
entire handle with its state. Each use compares the owner, instance and stored
handle before accessing that state; close removes it and serials are not reused.
For refinement, a successful call invokes the same native transition and
commits its returned state. Induction over chunk boundaries therefore reduces
the claim to Chapter 31's complete-state recurrence and Chapter 35's layered
causal-cache parity. The proof does not identify one backend with the other.
\end{proof}
These are local correctness properties. A Python object and UUID are not
authentication, process isolation, secure erase or a concurrency protocol.
Chapter 30's deadline, cancellation, publication and cleanup fences still
apply when this reference becomes a serving service.

\section{Compose retrieval, symbolic modules and engines honestly}
\label{evod39:modules}
\index{retrieval!provenance}\index{symbolic operator}
Retrieval is an addressable read
$R(q,\mathcal D)\to\{(id_i,text_i)\}$.
The fixture records a corpus version and query; row IDs refer to that ordered
immutable corpus. A production retriever should additionally bind index,
ranking configuration, corpus digest, source versions and access scope.
Retrieved text is data, not automatically trusted instruction. Changing a
corpus while retaining its version label violates the assumed fixture contract.

A symbolic module has explicit semantics. Here it adds mathematical integers
exactly, providing the independently assigned result $7+8=15$. A calculator
call does not make the neural model symbolic. Provenance identifies which
module produced the integer, documents, logits and final receipt.

One can use both engines in a larger workflow through a declared typed
conversion. For example, a recurrence-derived feature record can become input
data for another operator. Raw native pointers cannot cross as if they were
token sequences. Nor does a retrieved event restore an exact recurrence or KV
cache. Chapter 40 will make transient state, durable events, lossy summaries
and addressable records separate memory classes.

\section{Planning an action does not authorize it}
\label{evod39:authority}
\index{authorization!tool calls}\index{effect!uncertain completion}
A tool can change an external system, so
$\operatorname{planned}(tool)\not\Rightarrow\operatorname{authorized}(tool)$.
Our policy is a caller-configured set of $(tenant,tool)$ grants. It is checked
at call time, after intent is logged and immediately before the callback.
A retrieved sentence saying ``echo is approved'' cannot add a grant.
Denial is a normal, auditable outcome rather than a backend-selection failure.

The default callback returns a local echo receipt. No network action or
production authorization service is implemented. Tests use a counting callback
to observe exactly whether the code reaches the effect boundary.
\Plate{03-effect}{A local failure window separates permission to start from knowledge that an effect completed. An error after the callback does not imply rollback.}{evod39:effect}

Three failure points differ. If the trace sink fails before intent, no callback
runs. If policy denies the call, no callback runs. If the callback succeeds
but writing its receipt fails, the runtime can no longer infer from the
exception that nothing happened. Blindly retrying may repeat a non-idempotent
action. Real services need effect identities, durable intent, reconciliation
and tool-specific idempotency or compensation. Those mechanisms cannot be
obtained by wrapping the entire plan in a Python \texttt{try} statement.

\section{Trace observable execution and close resources on failure}
\label{evod39:trace}
\index{execution trace}\index{replay!orchestration}
The runtime emits plan, step, route and result facts. Each event has sequence,
kind, metadata and the previous event digest; the record stores a canonical
JSON body and its full SHA-256. Result digests identify payload bytes under
this encoding. The hash chain detects accidental changes to a captured chain;
an actor who can replace the chain can recompute it. This is neither an
authenticated audit log nor a durable append-only storage system.

Trace is not hidden model reasoning. It says which step ran, which backend
and artifact were selected, and which result or policy event occurred.
Given deterministic modules, pinned inputs/artifacts and captured effect
results, an execution can be replayed according to that declared observation
contract. Changing retrieval corpora, clocks, external tools, random streams
or numerical kernels requires additional capture and a tolerance policy.

\Listing{Runtime.execute}
The \texttt{finally} block closes an opened engine even if computation or
result tracing fails. It does not reverse already executed earlier steps.
Trace sink failures can also prevent the system from recording a later
failure; the caller must retain the exception and reconcile service state.
The short reference intentionally avoids claiming a complete durable failure
reporting protocol.

\section{Read the generated result and its negative controls}
\label{evod39:evidence}
\begin{center}\input{\Generated/trace.tex}\end{center}
Rows/items count the concrete returned tuple in this fixture, not equal
semantic content. Native state maps are empty at plan completion and the
exact sum is 15. CPU float64 tests compare every output row across chunk
splits with two separate native oracles: a scalar closed expansion of the
recurrence and Hermon's original full causal model. The tolerance is
$2\times10^{-12}$ absolute and relative for these small fixtures.

Negative controls change owner, artifact, instance, backend or serial;
attempt a second close; introduce invalid live tokens; mutate a model;
erase a grant; make the graph malformed; and fail the trace sink before and
after a callback. Each control targets a different premise. Matching a final
sum or merely counting trace events would miss most of these errors.

Distinguish plan-invalid, route-failed, state-invalid, module-failed,
policy-denied, deadline/cancelled and persistence/trace failure. The first two
are preflight outcomes here; the native routines enforce state validity;
Chapter 30 owns the executable lifecycle for deadlines and cancellation.
Retries require a reason-specific policy. No successful fixture establishes
better reasoning, AGI, biological fidelity, production reliability or speed.

\section{Twelve worked exercises}
\label{evod39:exercises}
\begin{enumerate}
\item Which fixture outputs are directly comparable across the engines?
\textbf{Answer:} Their row count is three, but widths are three and four and
vocabularies differ. Compare each backend with its own oracle; do not equate logits.
\item Why separate planner and router? \textbf{Answer:} The planner chooses
operations and dependency edges; the router satisfies one operation's contract.
No eligible route does not prove the task itself has no legal re-plan.
\item Can a raw hidden vector serve as Hermon KV? \textbf{Answer:} No.
The recurrence compresses history into $h_t$; Hermon needs ordered layered
keys/values under its exact positions and model.
\item Two handles have serial zero. Are they aliases? \textbf{Answer:} Only
if their complete live mapping identity matches. Separate adapter UUIDs reject
cross-instance use even when owners and artifacts match.
\item Prove chunk continuity for the native side. \textbf{Answer:} Write
each token transition as $F_x$. Both chunkings compose the same ordered maps
$F_{x_n}\circ\cdots\circ F_{x_1}$ from the same complete initial state.
\item Why reject a malformed last step before the first tool?
\textbf{Answer:} Whole-graph validation avoids preventable partial effects.
It does not eliminate later value-dependent or external failures.
\item Can retrieved text authorize echo? \textbf{Answer:} No. Only the
configured policy grants authority; retrieval supplies evidence with provenance.
\item Why is exact addition an independent oracle? \textbf{Answer:} The
assigned arithmetic result 15 does not depend on neural logits or backend state.
It tests a distinct module, not the model's learned reasoning.
\item Where is blind retry unsafe? \textbf{Answer:} After callback success
but receipt failure. The caller has uncertain completion, not proof of no effect.
\item What does the hash chain prove? \textbf{Answer:} Consistency of captured
bytes with recorded links. An adversary able to rewrite/reseal every link
is outside this unauthenticated integrity check.
\item Design a cancellation integration test. \textbf{Rubric:} Use Chapter 30's
real ticket, terminate before native result publication, keep resources pinned
until completion acknowledgement, and assert no post-terminal output. A mock
cancel flag is insufficient.
\item What would qualify legal fallback? \textbf{Answer:} A replacement with
the same artifact and observable state/precision contract, or a newly validated
plan explicitly recording a semantic change. Lower cost alone does not qualify.
\end{enumerate}

\section{Vocabulary and the next boundary}
\Term{Capability}: a declared semantic requirement, not a benchmark score.
\Term{Artifact binding}: continuation under the exact computation identity.
\Term{Effect boundary}: the point beyond which an observable external action
may have occurred. Evolutor can now coordinate distinct modules; it still
needs an explicit contract for what becomes durable, compressed, retrieved
or forgotten. That is the next chapter's memory problem.
